\section{谱距离估计与可靠上界的数值结果}
\label{sec:results}

本节报告一个固定候选的谱距离可靠上界，不把特殊算例改写成一般模型。算例取
\begin{align*}
  \beta&=\frac12,
  &d&=L=1,
  &\rho_{\mathrm e}&=1,
  &\rho_{\mathrm{disk}}&=17.
\end{align*}
左右周期胞元各含一个圆形介质柱；输入中半径的精确 dyadic 值显示为 $0.2$。中心胞元为空，即
$\Omega_0=\varnothing$。第~\ref{sec:error}~节的定理允许一般满足法向邻域假设的单个
$C^2$ 介质柱，中心也允许含一个形状不同的 $\Omega_0$；下面的解析尾界、圆周参数化和
空中心质量下界则是本算例专用的可靠上界计算步骤，不能直接推广到任意单连通形状。

固定的 binary64 候选按其精确值解释为
\begin{align*}
  \widehat k
  &=1.83277028910815698026226527872495353221893310546875,\\
  \mu=\widehat k^2
  &=3.3590469326375973209569763937462493\ldots.
\end{align*}
中心胞元宽度为 $2a=1$。几何法向邻域半宽的精确 dyadic 值显示为
$\delta_{\mathrm{col}}=0.04$，参考长度为 $\ell_*=1$。候选计算交付固定的原始密度样本，
可靠上界计算阶段另用 280 个目标求积节点评价连续材料跳跃；原始场离散与后续求积承担不同任务。
材料局部整函数级数的实现截断阶数
\texttt{Q} 为 10，即以阶数 10 作为该局部解析展开的截断参数；复平面条带宽度
参数为 $1/5$，Ewald 的 split、reciprocal、image、local 与 Taylor 截断依次为
$2,5,3,18,18$，区间运算精度为 192 bit。这里的代码变量 \texttt{Q} 只是级数阶数，不是
式~\eqref{eq:trace-defect-product} 的 $Q$ 或下文的 $\overline Q$。人工竖直界面保留
$-256\leq m\leq256$ 的 Rayleigh 模态；左右各显式计算第 $0$ 至第 $23$ 个周期胞元，
显式胞元数 $24$ 由预先给定的序列 $16,24,32,48,64$ 选择，传播尾的矩阵分块长度为 $1$，
其余无限胞元由完整 $P_\pm$ 的收缩上界覆盖。局部与准周期 Green 函数的各项、传播尾、质量和最终公式均使用
Arb 的 midpoint--radius 球算术向外包围；该算术框架见 Johansson
\cite{Johansson2017Arb}，但本文定理及针对当前核的误差分解并非该文献的结论。

\paragraph{备注（密度样本数）.}
在当前圆形空中心算例的候选计算阶段，分辨率测试表明
$N_\eta=256$ 在重构稳定性与计算成本之间给出较好的经验折衷，因此本次计算固定这一
样本数，并按式~\eqref{eq:exact-density-dft} 构造有限三角密度。这个数值选择不是
第~\ref{sec:error}~节定理的假设，也不替代可靠上界计算阶段独立的目标求积与误差包围。

几何条件的验证同样使用上述精确 dyadic 输入。圆的曲率上界取
$\kappa_*=1/R$，覆盖重数为 $N_{\mathrm{ov}}=1$，并以区间运算验证
$\vartheta_{\mathrm{col}}=1-\delta_{\mathrm{col}}/R>0$；这里只保留符号关系和严格正性，
不把相应十进制显示值当作未经包围的精确等式。计算中还分别核验了材料曲线与非零
$x$ 向周期副本、上下 $y$ 向周期副本、材料法向邻域与人工竖直界面法向邻域，以及不同
人工竖直界面法向邻域之间的分离性。因此命题~\ref{prop:positive-shift-trace-bound} 的共同
法向邻域和有限覆盖假设在本例中得到验证。完整传播矩阵 $P_\pm$ 的可靠收缩界既控制所有未
显式计算胞元的界面跳跃尾，也保证原始场的胞元质量平方可和，从而核验
$u_{h,M}\in\mathcal H$。

最终可靠单侧输入与组合结果为
\begin{align}
  \overline C
  &\leq2.34995201038114,
  &\overline\varepsilon_\partial
  &\leq1.30332365404084\times10^{-8},
  \notag\\
  \underline U
  &\geq1.24751965127352,
  &\overline Q
  &\leq3.06274804099054\times10^{-8},
  \notag\\
  \underline D
  &\geq1.24751962064604>0,
  &R_{\mathrm{spec}}
  &\leq1.64933909531673\times10^{-7}.
  \label{eq:v3-certified-summary}
\end{align}
全部便读上界均向上舍入、下界均向下舍入；实际组合使用完整 dyadic 端点，而不是上式的
十进制截断。正分母条件已经通过验证。相应向外区间为
\begin{align}
  [3.359046767703687789,\;3.359047097571506853],
  \label{eq:v3-certified-interval}
\end{align}
它与 $\sigma(A(\beta))$ 相交。由于本次计算没有另行证明本质谱隙，式
~\eqref{eq:v3-certified-interval} 不能单独证明其中存在离散导模特征值，更不能证明唯一性、
编号或谱隙。

界面缺陷平方上界的四个合计分项为
\begin{center}
  \begin{tabular}{lr}
    \hline
    分项 & 可靠平方上界\\
    \hline
    材料界面值跳跃
      & $2.86590056317554\times10^{-18}$\\
    材料界面普通法向导数跳跃
      & $1.66900393921493\times10^{-16}$\\
    人工界面值跳跃
      & $5.14806044930848\times10^{-31}$\\
    人工界面普通法向导数跳跃
      & $9.89602335658444\times10^{-20}$\\
    \hline
  \end{tabular}
\end{center}
材料分项直接在圆周上作有符号的内外迹组合和连续曲线求积，没有对目标材料跳跃再作
傅里叶截断；准周期 Green 函数本身的 Rayleigh 表示、源密度离散 Fourier 变换，以及人工
界面的 Fourier 与 Rayleigh 尾仍完整计入。当前界的主要保守性来自阶数 10 的材料普通法向
局部连续级数尾，其次包括解析目标求积余项和完整传播尾；不能把可靠上界与先前估计的差异
全部归因于浮点舍入。

作为对照，早期 v2 计算采用普通双精度、1026 个材料节点、局部级数阶数 32、左右各显式
16 个胞元、长度 32 的传播分块和 $|m|\leq256$ 的人工界面模态。其公式评价为
\begin{align*}
  R_{\mathrm{spec}}
  &=4.0547573451465766\times10^{-9}.
\end{align*}
相应名义区间为
\begin{align*}
  [3.3590469285828397,\;3.3590469366923545].
\end{align*}
该结果没有对特殊函数、传播尾、舍入方向和质量下界作完整单侧包围，因而不是可靠上界。
v3 为用球算术计算可靠上界而采用了不同的节点数、局部阶数和尾界配置；它的区间较宽反映的是
完整误差所有权和保守包围，并不表示候选精度下降，也不使 v2 的估计自动成为可靠上界。

最终验证记录为
\texttt{test/i6/resolvent-v3/output/attempt-1/run-10/certificate.json}；独立标量审计文件
\texttt{scalar-audit.json} 对 96 个球记录的精确端点和组合公式全部通过。复现入口为同目录的
\texttt{run\_with\_limits\_v3.py}，固定输入、源码快照、依赖版本和资源记录均随 run-10
保存；本次环境记录为 MATLAB R2023b、Python 3.14 与 python-flint 0.9.0。该可靠上界依赖
Arb/FLINT 的包含性算术约定和已审查的公式实现，并不是整个软件栈的形式化
验证；它也不为此前的候选搜索过程提供正确性保证。
